\documentclass[11pt,a4paper]{article}
\usepackage[utf8]{inputenc}
\usepackage[T1]{fontenc}
\usepackage{amsmath,amssymb}
\usepackage{booktabs}
\usepackage{hyperref}
\usepackage[margin=1in]{geometry}

\title{Chromavox: Unblended Color-Set Superposition Outperforms Token Sequences on a Disentanglement Task}
\author{Sean Kellogg}
\date{}

\begin{document}
\maketitle

\begin{abstract}
We introduce \textbf{chromavox} (chroma + voxel), a colored voxel semantic-space representation in which each cell of a 3D grid carries a bounded SET of color entries (up to $K{=}2$), each with its own hue, brightness, and alpha. Colors in a set coexist without blending---red AND blue, never purple---matching the ``simultaneous color experience'' reported in the synesthesia literature and providing a native representation for uncertainty and superposition.

We test three claims in sequence: (1) the chromavox space can hold semantic structure projected from a pretrained text embedding (V0, structure preservation $\rho{=}0.984$); (2) a learned MLP encoder can produce chromavox fields from 64-dim matryoshka embeddings while preserving 95--97\% of pairwise distance structure (V1, $\rho{=}0.954$--$0.973$); and (3) on a superposition task where two streams share the same spatial positions, the chromavox representation decisively outperforms a param-matched token baseline (C1a, relative error $0.114$ vs $1.004$, token exact-match $0\%$).

The superposition result is the first demonstration that an unblended color-SET representation offers a task-specific advantage over flat token sequences. We honestly report two negative results on the same substrate (differentiation and integration, where tokens dominate) and discuss the scope limitations of the positive finding.
\end{abstract}

\section{Introduction}

Representation determines capability. Abacus embeddings~\cite{abacus} showed that the right positional embedding can take a transformer from near-0\% to near-perfect on arithmetic. Lample \& Charton~\cite{lample2020} beat Mathematica at symbolic integration through a representation choice (prefix-notation tokenization). The geometry of the representation, not the model architecture, was the lever.

This paper asks whether a \textit{structured} representation---hue as an ordinal axis, color-sets as un-collapsed superpositions, brightness as salience, an ego-centered spatial grid---gives a model capabilities that a plain token sequence does not. The representation is motivated by synesthesia: the author has synesthetic perception in which multiple colors coexist at the same spatial location without blending. We treat this as a computational primitive (a bounded SET, not an average) and test whether it earns its keep.

\section{The Chromavox Representation}

A chromavox field is an $8^3$ grid (512 cells). Each cell carries:
\begin{itemize}
\item \textbf{Density:} float $[0,1]$, quantized to 2 significant figures.
\item \textbf{Color-set:} up to $K{=}2$ entries, each (hue\_cos, hue\_sin, brightness, alpha).
\item \textbf{Empty cells:} density$=0$, no color entries.
\end{itemize}

Key properties:
\begin{itemize}
\item Colors never blend. Two entries in the same cell coexist as distinct values---the set $\{\text{red}, \text{blue}\}$ is not purple. This is the structural claim that distinguishes chromavox from additive embedding methods.
\item Spatial positions are discrete (512 cells). Thoughts occupy specific regions; simple thoughts use few cells, complex ones pack the space.
\item Quantization to 2 sig figs provides a finite, bounded representation.
\end{itemize}

Effective dimensionality: $512 \times 9 = 4{,}608$. But the constraints (sigmoid bounding, cos/sin hue, 2-sig-fig quantization) reduce effective capacity below the raw dimensionality. The question is how much.

\section{Experiments}

All experiments use pre-registered gates, param-matched controls, and deterministic seeds (seed 1337).

\subsection{V0: Can chromavox space hold semantic structure?}

\textbf{Method.} Embed 2,000 STS-B test sentences~\cite{stsb} with nomic-embed-text v1.5 (matryoshka-truncated to 64 dims). Project into chromavox via a fixed orthogonal projection (Johnson--Lindenstrauss), with constraints (sigmoid, cos/sin, 2-sig-fig quantization) baked into the forward pass. Measure Spearman $\rho$ of pairwise distances: projected vs original embedding.

\textbf{Arms.} A (voxel-full), B (density-only), C (unconstrained linear), D (random projection), E (identity/no projection), F (no quantization).

\begin{table}[h]
\centering
\begin{tabular}{lccc}
\toprule
Arm & $\rho$ (vs emb) & $\rho$ (vs STS-B) & dims \\
\midrule
A (voxel-full) & 0.984 & $-$0.824 & 4,608 \\
B (density-only) & 0.999 & $-$0.824 & 4,608 \\
C (unconstrained) & 1.000 & $-$0.824 & 4,608 \\
D (random) & 0.031 & $+$0.051 & 4,608 \\
E (identity) & 1.000 & $-$0.824 & 64 \\
F (no quant) & 0.984 & $-$0.824 & 4,608 \\
\bottomrule
\end{tabular}
\end{table}

\textbf{Finding.} Chromavox constraints preserve 98.4\% of embedding structure ($\rho{=}0.984$ vs random $0.031$). Quantization costs $<0.001$ $\rho$. The space is not broken for semantics. However, color channels do not add capacity beyond density alone (0.984 vs 0.999)---512 density dimensions already exceed the 64-dim input. This is a capacity test, not an emergence test: any sufficiently high-dimensional space would pass. It establishes the prerequisite, not the claim.

\subsection{V1: Can a model learn to produce chromavox fields?}

\textbf{Method.} Train an MLP encoder ($64 \to 256 \to 4{,}608$, ${\sim}131$K params) on 5,749 STS-B training pairs with two objectives: (A) distance-preservation loss (minimize MSE of pairwise distances) and (B) autoencoder reconstruction loss (encode then decode back to 64-dim). Constraints (sigmoid, cos/sin, straight-through quantization) baked into the encoder forward pass. Evaluate on 2,552 held-out test sentences.

\begin{table}[h]
\centering
\begin{tabular}{lccc}
\toprule
Arm & $\rho$ (vs emb) & dims & params \\
\midrule
A (learned-dist) & 0.954 & 4,608 & 131K \\
B (learned-autoenc) & 0.973 & 4,608 & 262K \\
C (fixed-ortho) & 0.983 & 4,608 & 0 \\
D (random) & 0.002 & 4,608 & 0 \\
E (identity) & 1.000 & 64 & 0 \\
\bottomrule
\end{tabular}
\end{table}

\textbf{Finding.} A learned MLP encoder preserves 95.4\% of embedding structure (arm A) and 97.3\% via autoencoder reconstruction (arm B), reaching 97\% of the fixed-orthogonal ceiling (arm C: 0.983). All gates pass. The autoencoder objective (B) outperforms the distance-preservation objective (A), suggesting that reconstructable information content is a stronger training signal than pairwise distance matching. A model can learn to fill chromavox space with meaning.

\subsection{C1a: Superposition---chromavox beats tokens at disentanglement}

\textbf{Method.} Two functions $f, g$ sampled on the same 32 x-bins, voxelized into the SAME chromavox cells. Each cell holds two color-set entries (one for $f$, one for $g$). Task: output the derivative of $f$ ONLY. The model must attend to the $f$-stream and ignore the $g$-stream. Param-matched token baseline (10.67M vs 10.90M, within 2.2\%). Token baseline interleaves $[f_0, g_0, f_1, g_1, \ldots]$ as a flat sequence.

Three token serialization controls: interleaved (original), concatenated $[f_0, \ldots, f_{31}, g_0, \ldots, g_{31}]$, and a corrected-color (cos/sin hue) voxel arm. An energy/wavelength encoding (200--1000nm, UV--IR) is also tested.

\begin{table}[h]
\centering
\begin{tabular}{lcccc}
\toprule
Arm & rel.\ err & exact & extrap & params \\
\midrule
voxel (cos/sin) & 0.114 & 0.889 & 0.563 & 10,674,065 \\
voxel (energy) & 0.116 & 0.887 & 0.565 & 10,674,065 \\
token (interleaved) & 1.004 & 0.000 & 1.004 & 10,903,115 \\
token (concatenated) & 1.004 & 0.000 & 1.004 & 10,903,115 \\
\bottomrule
\end{tabular}
\end{table}

\textbf{Finding.} The chromavox voxel arm achieves relative error 0.114 while BOTH token baselines completely fail (rel.\ err 1.004, exact-match 0.000). The token failure is not a serialization artifact---concatenated tokens fail identically to interleaved. The color-SET representation genuinely outperforms tokens at keeping co-located streams separate and extracting the correct one.

The energy/wavelength encoding (200--1000nm, UV--IR) ties the cos/sin encoding (0.116 vs 0.114), confirming the advantage comes from the SET structure, not the specific color values.

\subsection{Negative results (honestly reported)}

Two other operators on the same substrate went to tokens:

\begin{itemize}
\item \textbf{Differentiation (C0):} token baseline 0.006 vs best voxel 0.10. Tokens are $15\times$ better at pointwise differentiation. The voxel substrate offers no advantage on local operators where attention already has full context.
\item \textbf{Integration (C1b):} token baseline 0.007 vs voxel 0.225. The design hypothesis that integration (nonlocal accumulation) would favor the substrate was falsified by the data.
\end{itemize}

We also tested whether meaning-locked invariance pairs (same-clip AudioCaps captions) create more semantic structure than random pairs (B1). They do not (same $\rho$ 0.607 vs shuffled 0.625, delta $-$0.017). The structure came from the pretrained encoder (CLAP), not from the pair mechanism. The binding hypothesis remains unproven.

\section{Discussion}

The positive result is narrow but genuine: on a task that exercises chromavox's one distinctive capability (multiple things in one cell, kept separate), it decisively outperforms a param-matched token baseline that completely fails. This is the first demonstration of a task-specific advantage for an unblended color-SET representation.

The negative results define the boundary: the advantage does not extend to local pointwise operators (differentiation, integration) where token attention already excels. The substrate is not a general-purpose computational engine---it is a representation with a specific structural capability (superposition/disentanglement) that tokens lack.

The V0 and V1 results establish that chromavox space can hold and recover semantic structure from pretrained embeddings, with a learned encoder preserving 95--97\% of pairwise distance structure through the constrained projection. This is necessary but not sufficient: it proves the space is not broken for semantics, not that its structure is better than a plain high-dimensional vector.

\section{Scope Limitations}

\begin{enumerate}
\item The superposition result is on synthetic calculus (two functions sharing 32 bins), not semantic content. Whether the advantage extends to semantic ambiguity (multiple word senses, candidate referents) is the open question.
\item V0 and V1 use a frozen pretrained encoder (nomic-embed-text v1.5) as input. Whether a model trained end-to-end (text $\to$ chromavox, no frozen encoder) can produce meaningful fields is untested.
\item No dynamic belief-update mechanism has been tested. The vision of ``hold possibilities, update with new info, converge'' requires temporal modeling that none of these experiments exercise.
\item STS-B is a sentence-similarity benchmark, not a reasoning or downstream-task benchmark. Structure preservation on STS-B does not imply task performance.
\end{enumerate}

\section{Related Work}

\begin{itemize}
\item \textbf{Abacus embeddings}~\cite{abacus}: representation geometry determines arithmetic capability. Chromavox tests the same principle for a different representation.
\item \textbf{Lample \& Charton}~\cite{lample2020}: prefix-notation tokenization beats Mathematica at symbolic integration. The token baseline in C1a is in this lineage.
\item \textbf{Anthropic toy models of superposition}~\cite{toy}: features superpose in neural networks as directions in activation space. Chromavox makes superposition an explicit, bounded SET rather than an emergent direction.
\item \textbf{Ramachandran \& Hubbard}~\cite{ramachandran}: synesthesia as cross-modal binding. Chromavox is motivated by but does not claim to model synesthetic perception.
\item \textbf{Eagleman et al.}~\cite{eagleman}: clinical diagnostic battery for synesthesia (consistency, automaticity). The operational definition of ``machine synesthesia-like binding'' in this project draws from this framework.
\end{itemize}

\section{Reproducibility}

All experiments use deterministic seeds (seed 1337) and are reproducible on
a standard laptop with Python 3.12, PyTorch, NumPy, and SciPy. V0 and V1
run on CPU in under 30 minutes total. C1 runs on a single GPU
(AMD Radeon 8050S iGPU, ${\sim}2$h). Complete source code, pre-registered
designs, per-epoch logs, and machine-readable result files are available
at the author's repository.\footnote{Source code and machine-readable
results are available at \texttt{github.com/Ziggibot0/synesthetic-ai}
(commit 4881fc1). As an evolving research repository, the code may
diverge from the state described in this paper over time; the cited
commit hash is the authoritative reference for reproducibility.} Key implementation details:

\begin{itemize}
\item \textbf{V0:} 2,000 STS-B test sentences, nomic-embed-text v1.5
(64-dim matryoshka), fixed orthogonal projection, 2-sig-fig quantization.
${\sim}5$ min CPU.
\item \textbf{V1:} 10,566 STS-B train sentences, MLP encoder
($64 \to 256 \to 4{,}608$), 200 epochs, AdamW lr$=3\times10^{-4}$.
${\sim}20$ min CPU.
\item \textbf{C1a:} 20,000 training samples (sympy-generated polynomials
and trigonometric functions, 32 x-bins), 6-layer transformer encoder
($H{=}384$, FFN$=1{,}536$, 8 heads), 60 epochs, batch 128.
Param-matched: voxel 10,674,065 vs token 10,903,115 (within 2.2\%).
\end{itemize}

\bibliographystyle{plain}
\begin{thebibliography}{6}

\bibitem{abacus}
S.~Jelassi et~al., ``Transformers Can Do Arithmetic with the Right Embeddings,''
arXiv:2405.17399, 2024.

\bibitem{lample2020}
G.~Lample and F.~Charton, ``Deep Learning for Symbolic Mathematics,''
in \textit{Proc.\ ICLR}, 2020.

\bibitem{toy}
Anthropic, ``Toy Models of Superposition,''
transformer-circuits.pub, 2022.

\bibitem{ramachandran}
V.~S.~Ramachandran and E.~M.~Hubbard, ``Synaesthesia---A Window into Perception, Thought and Language,''
\textit{J.\ Consciousness Studies}, vol.~8, no.~12, 2001.

\bibitem{eagleman}
D.~M.~Eagleman, A.~Kagan, S.~Nagarajan, S.~Sagaram, and A.~Sarma, ``A standardized test battery for the study of synesthesia,''
\textit{J.\ Neuroscience Methods}, vol.~159, no.~1, 2007.

\bibitem{stsb}
D.~Cer et~al., ``Semantic Textual Similarity Benchmark,''
in \textit{Proc.\ ACL}, 2017.

\end{thebibliography}

\end{document}